% Image credits for the photographs and AI illustrations of Book 2
% (University Chemistry, Year 1). Machine-readable license records live in
% images/book2/CREDITS.md; keep the two files in sync. The Book 2 writer
% owns this file: add one \item per photograph (in an itemize, after the
% first paragraph) as photographs land.
\chapter*{Image Credits}

The drawings in this book are original. Photographs, where used, are used
under their respective free licenses, with thanks to their authors; full
source links and license URLs for every photograph are listed in
\texttt{images/book2/CREDITS.md} in the book's repository.
The hazard pictograms are those of the Globally Harmonized System of
Classification and Labelling of Chemicals, as drawn for the
\texttt{ghsystem} package of \TeX\ Live.

\bigskip
Alongside the photographs and the original drawings, some figures are
AI-generated illustrations, produced with OpenAI image generation from
prompts written by the book's editors, and reviewed for chemical
accuracy before inclusion. The full prompt for every generated image is
recorded in \texttt{images/book2/ai/PROMPTS.md} in the repository.

\bigskip
\noindent Photographs, by chapter:
\begin{itemize}\small
  \item Ch.~1: Niels Bohr, about 1922, AB Lagrelius \& Westphal, public domain
        (Wikimedia Commons).
  \item Ch.~2: Linus Pauling, 1962, Nobel Foundation, public domain
        (Wikimedia Commons).
  \item Ch.~5: native copper specimen, by Flipin, OG, CC0 (Wikimedia
        Commons).
  \item Ch.~6: halite crystals, fluorite and an octahedral diamond, by
        James St. John, CC BY 2.0; graphite, by Zbynek Burival,
        CC BY-SA 4.0; snow crystals by Wilson A. Bentley (1902), public
        domain (all Wikimedia Commons).
  \item Ch.~8: Svante Arrhenius, 1909, Meisenbach Riffarth \& Co., public
        domain (Wikimedia Commons).
  \item Ch.~11: silver chloride, bromide and iodide precipitates, by
        Rrausch1974, CC BY-SA 3.0 (Wikimedia Commons).
  \item Ch.~12: copper(II) hydroxide suspension, by Alvy16, CC BY 4.0;
        copper ammine solution, by Chemicalinterest, public domain (both
        Wikimedia Commons).
  \item Ch.~13: Walther Nernst, 1889, Smithsonian Libraries, public domain
        (Wikimedia Commons).
  \item Ch.~16: Jacobus Henricus van 't Hoff, photograph published 1911,
        public domain (Wikimedia Commons).
  \item Ch.~21: Victor Grignard, public domain (Wikimedia Commons).
  \item Ch.~25: Vladimir Markovnikov, portrait from an obituary notice
        published in 1905, public domain (Wikimedia Commons).
  \item Ch.~26: lithium ponds of the Salar de Atacama, NASA Earth
        Observatory, public domain; sodium in mineral oil, by Justin Urgitis,
        CC BY-SA 2.5; flame colours, by Hegelrast, CC BY-SA 4.0 (all
        Wikimedia Commons).
  \item Ch.~27: pisolitic bauxite, by James St. John, CC BY 2.0; borax, by
        Aram Dulyan, public domain; quartz, by Stephanie Clifford, CC BY 2.0;
        white phosphorus, Hi-Res Images of Chemical Elements, CC BY 3.0;
        Fritz Haber, The Nobel Foundation, public domain (all Wikimedia
        Commons).
  \item Ch.~28: sulfur mining in the Kawah Ijen crater, by S\'emhur,
        CC BY-SA 4.0; chlorine in an ampoule, by W. Oelen, CC BY-SA 3.0;
        bromine in an ampoule, by Jurii, CC BY 3.0; iodine crystals, by
        Dnn87, CC BY 3.0 (all Wikimedia Commons).
  \item Ch.~29: heated bench (Kofler bench), by HBR, CC BY-SA 3.0; Abbe
        refractometer, by Foreade, CC BY-SA 4.0 (both Wikimedia Commons).
\end{itemize}
\clearpage
